\section{An explicit reduction to 853 variables}
\label{sec:explicit-853}

Two further elementary steps combine the private anchor root with the
Lucas root. Their proof concerns rational solutions. It is not a claim
about geometric fibers over an algebraic closure.

\subsection{The pole denominators are automatic over the rationals}

\begin{lemma}\label{lem:pole-denominator-automatic}
After clearing the denominator $\delta=1+bt^2$ from
\eqref{eq:circuit-pole}, its two equations already imply
$\delta\ne0$ over $\Q$.
\end{lemma}
\begin{proof}
Put $B_p=17b$ and
$A=B_p\delta^2-b^2(1-bt^2)^2$. The cleared second equation is
\[
 B_pq^2\delta^4
 =(h_p^2\delta^2-B_pe_pA)(h_p^2\delta^2+3B_pe_pA).
\]
If $\delta=0$, then $bt^2=-1$, so $b,t,B_p$ are nonzero and
$A=-4b^2$. The equation becomes $0=-48B_p^2b^4e_p^2$, whence
$e_p=0$. The first equation now gives $p^2=17b=-17/t^2$,
or $(pt)^2=-17$, impossible in $\Q$.

This also applies to the reduced rational numerators used by the
circuit. In each actual second conic, the denominator contains the
numerator of $\delta$ to exactly the fourth power, and every other
factor is among the retained nonzero conditions. The first-conic
denominator uses only retained factors. Multiplying back yields the
displayed polynomial identity even when $\delta=0$. These exact
denominator-factor statements are checked in the independent verification.
\end{proof}

More generally the same proof works with $B_p=cb$, over a field of
characteristic different from two and three, whenever $-c$ is a
nonsquare. It deliberately does not give the corresponding assertion
over every number field.

Remove precisely the $42$ pole-denominator conditions from the product
$G$ of the preceding section, and call the remaining product $G_0$.
Its factors depend only on the input, point coordinates, coordinate
denominator, and contact coordinates. In particular $G_0$ is independent
of all three proxies, all pole coordinates, $\theta$, and the private
anchor coordinates. This follows from the displayed construction and
was checked on the entire circuit. All other rational denominators
are covered by $G_0$; the omitted ones are enforced by
Lemma~\ref{lem:pole-denominator-automatic}.

\subsection{Combining two square roots and the nonzero conditions}

Write the Lucas equation as
\[
 \tau^2-a\tau+b_0=0,\quad
 a=s(s-2)z,\quad b_0=(s-2)^2z^2-1,\quad
 \Delta=a^2-4b_0.
\]
Thus $v=2\tau-a$ satisfies $v^2=\Delta$. The scaled anchor witness
satisfies $t_a^2=-N$, where $N$ is independent of $\tau$.
Choose its nonzero representative as in the preceding section.
Its sign can be selected so that $W=t_a+v\ne0$: the two choices
$t_a+v$ and $-t_a+v$ cannot both vanish when $t_a\ne0$.
Consequently no new finite selector is necessary.

The two square-root equations have the exact single relation
\begin{equation}\label{eq:joint-quartic}
 h(W)=W^4+2(N-\Delta)W^2+(N+\Delta)^2=0,
\end{equation}
with reconstruction
\[
 v=\frac{W^2+N+\Delta}{2W},\qquad
 t_a=\frac{W^2-N-\Delta}{2W},\qquad
 \tau=\frac{a+v}{2}.
\]
Indeed $4W^2(v^2-\Delta)=h(W)=4W^2(t_a^2+N)$.
Conversely, $W=t_a+v$ and the two square equations imply
\eqref{eq:joint-quartic} and these reconstruction formulas.

Set $W=1/(G_0U)$ and define
\[
 \Psi_*=1+2(N-\Delta)G_0^2U^2+(N+\Delta)^2G_0^4U^4,
 \qquad
 \tau=\frac{1+2aG_0U+(N+\Delta)G_0^2U^2}{4G_0U}.
\]
The equation $\Psi_*=0$ forces $G_0U\ne0$ and is the reciprocal
form of \eqref{eq:joint-quartic}. Substitute the displayed rational
formula for $\tau$ into all the other residuals. Choose a syntactic
upper bound $k$ for each residual's degree in $\tau$, and multiply
its substituted expression by $(4G_0U)^k$ to obtain a polynomial
numerator. At addition gates take the maximum of the children's
bounds; at multiplication gates add them; at power gates multiply
by the exponent. This specifies the emitted circuit even if
cancellation lowers the actual degree. The operation is reversible
on $G_0U\ne0$.
Delete the old Lucas residual and the anchor residual; denote the
remaining $295$ polynomial numerators by $E'_i$. The branch polynomial
and the full polynomial are
\[
 F^*_{j,c}=\sum_{i=1}^{295}(E'_i)^2+\Psi_*^2,
 \qquad
 P_{853}=\prod_{j=0}^{211}\prod_{c\in\mathcal C}F^*_{j,c}.
\]

\begin{proposition}\label{prop:853-explicit}
The integer polynomial $P_{853}$ has 853 total variables and the
same rational input projection as $\Src$. The supplied finite
arithmetic circuit has degree at most $894090701148160$ and
coefficient height at most $2^{29154205340326912}$.
\end{proposition}
\begin{proof}
The forward construction may change only the three private anchor
witnesses, making $t_a\ne0$ and then choosing its sign to make
$W\ne0$. Set $U=1/(G_0W)$; every transformed residual vanishes.
Conversely, $\Psi_*=0$ gives $G_0U\ne0$. Reconstruct $W,v,t_a$
and $\tau$ by the formulas above, then $u_a=t_a/S_0$. They satisfy
the old anchor and Lucas equations. Undoing each homogenization gives
the other original numerator equations. The retained guards and
Lemma~\ref{lem:pole-denominator-automatic} make every original rational
division valid. The source point-chart equivalences then apply as
before. These operations prove equality of rational input projections.

Two coordinates, $u_a$ and $\tau$, have been replaced by the single
coordinate $U$. The count is
\[
 633+210+2+2+1+1+2+1+1=853,
\]
respectively point coordinates, pole coordinates, two remaining anchor
coordinates, two remaining proxies, $\theta$, the coordinate denominator,
two contact coordinates, the input, and $U$. There are still $54272$
finite branches. Exact circuit bounds are propagated as in
Proposition~\ref{prop:854-explicit}; the degree and height statements
are upper bounds. The maximum homogenizing degree in $\tau$ is $1280$.
\end{proof}

The omission of the pole guards is justified here over $\Q$.
The 558-variable geometric-fiber argument retains those guards and
does not use this omission. This distinction prevents exceptional
geometric fibers from being silently reintroduced into that argument.
